\section{Introduction}
%P1: Background, importance of recommendation and CF
To alleviate information overload on the web, recommender system has been widely deployed to perform personalized information filtering~\cite{PinSage,YoutubeRS, PeterRec}.
The core of recommender system is to predict whether a user will interact with an item, e.g., click, rate, purchase, among other forms of interactions. As such, collaborative filtering (CF), which focuses on exploiting the past user-item interactions to achieve the prediction, remains to be a fundamental task towards effective personalized recommendation~\cite{NCF,VACF,NGCF,CMN}. 

%P2: Brief related work, matrix factorization -> neural nets. Introduce key idea of GCN and highlight it has been new SOTA. 
The most common paradigm for CF is to learn latent features (a.k.a. embedding) to represent a user and an item, and perform prediction based on the embedding vectors~\cite{NCF,DBLP:conf/www/ChengDZK18}.
Matrix factorization is an early such model, which directly projects the single ID of a user to her embedding~\cite{MF}.  Later on, several research find that augmenting user ID with the her interaction history as the input can improve the quality of embedding. For example, SVD++~\cite{SVD++} demonstrates the benefits of user interaction history in predicting user numerical ratings, and Neural Attentive Item Similarity (NAIS)~\cite{NAIS} differentiates the importance of items in the interaction history and shows improvements in predicting item ranking. 
In view of user-item interaction graph, these improvements can be seen as coming from using the subgraph structure of a user --- more specifically, her one-hop neighbors --- to improve the embedding learning. 

%P3: Argue some GCN designs are unnecessary for CF. Check SGCN paper how to argue it. The input to CF task is simple and less informative -- only ID to learn user/item embedding => may not useful to perform nonlinear feature transformation.
To deepen the use of subgraph structure with high-hop neighbors, Wang et al.~\cite{NGCF} recently proposes NGCF and achieves state-of-the-art performance for CF. It takes inspiration from the Graph Convolution Network (GCN)~\cite{GCN,GraphSAGE}, following the same propagation rule to refine embeddings: feature transformation, neighborhood aggregation, and nonlinear activation. Although NGCF has shown promising results, we argue that its designs are rather heavy and burdensome --- many operations are directly inherited from GCN without justification. As a result, they are not necessarily useful for the CF task. To be specific, GCN is originally proposed for node classification on attributed graph, where each node has rich attributes as input features; whereas in user-item interaction graph for CF, each node (user or item) is only described by a one-hot ID, which has no concrete semantics besides being an identifier. In such a case, given the ID embedding as the input, performing multiple layers of nonlinear feature transformation --- which is the key to the success of modern neural networks~\cite{ResNet} --- will bring no benefits, but negatively increases the difficulty for model training. 

%P4: Introduce our work
To validate our thoughts, we perform extensive ablation studies on NGCF. With rigorous controlled experiments (on the same data splits and evaluation protocol), we draw the conclusion that the two operations inherited from GCN --- feature transformation and nonlinear activation --- has no contribution on NGCF's effectiveness. Even more surprising, removing them leads to significant accuracy improvements. This reflects the issues of adding operations that are useless for the target task in graph neural network, which not only brings no benefits, but rather degrades model effectiveness. Motivated by these empirical findings, we present a new model named LightGCN, including the most essential component of GCN --- neighborhood aggregation --- for collaborative filtering. Specifically, after associating each user (item) with an ID embedding, we propagate the embeddings on the user-item interaction graph to refine them. We then combine the embeddings learned at different propagation layers with a weighted sum to obtain the final embedding for prediction. The whole model is simple and elegant, which not only is easier to train, but also achieves better empirical performance than NGCF and other state-of-the-art methods like Mult-VAE~\cite{VACF}. 

%P5: Summary of contributions.
To summarize, this work makes the following main contributions:
\begin{itemize}%[leftmargin=*]
    \item We empirically show that two common designs in GCN, feature transformation and nonlinear activation, have no positive effect on the effectiveness of collaborative filtering. 
    \item We propose LightGCN, which largely simplifies the model design by including only the most essential components in GCN for recommendation.
    \item We empirically compare LightGCN with NGCF by following the same setting and demonstrate substantial improvements. In-depth analyses are provided towards the rationality of LightGCN from both technical and empirical perspectives. 
    %XN: to revise
    %Further analyses show the advantages of LightGCN by being simple --- easier to train, less sensitive to hyper-parameters, and computationally more efficient. 
\end{itemize}

%SPACE
%The remainder of the paper is organized as follows. In Section~\ref{sec:preliminaries}, we provide empirical evidence on the unnecessarily heavy design of NGCF as the preliminaries. We then present our proposed LightGCN method and technically analyze its properties in Section~\ref{sec:method}. We conduct experiments to validate LightGCN in Section~\ref{sec:experiments}. Finally, we review related work in Section~\ref{sec:related} and conclude the paper in Section~\ref{sec:conclusion}. 